\clearpage
\VerdierPhysicalPageAnchor{0179}%
\begingroup\small\setlength{\parskip}{0pt}\setlength{\emergencystretch}{1.5em}
\setlength{\abovedisplayskip}{4pt plus 2pt minus 2pt}%
\setlength{\belowdisplayskip}{4pt plus 2pt minus 2pt}%

\VerdierSemanticAnchor{ch3d-1-2-4}
\noindent\textbf{1.2.4.}\quad
$D(\mathcal A)$ 的子范畴
$K^{,b}(\mathcal A)/\operatorname{Ac}(\mathcal A)$（相应地，
$K^{,+}(\mathcal A)/\operatorname{Ac}(\mathcal A)$，相应地，
$K^{,-}(\mathcal A)/\operatorname{Ac}(\mathcal A)$）记作
$D^b(\mathcal A)$（相应地，$D^+(\mathcal A)$，相应地，
$D^-(\mathcal A)$），并称为\emph{有界对象子范畴}（相应地，
\emph{下有界对象子范畴}，相应地，\emph{上有界对象子范畴}）。应当注意，
这些有界性条件仅施加于所考虑复形的上同调对象。

\VerdierSemanticAnchor{ch3d-1-2-5}
\medskip\noindent\textbf{1.2.5.}\quad
设 $X$ 和 $Y$ 是 $D(\mathcal A)$ 的两个对象。以
$\operatorname{Ext}_{\mathcal A}^{0}(X,Y)$ 表示群
$\operatorname{Hom}_{D(\mathcal A)}(X,Y)$。同样，以
$\operatorname{Ext}_{\mathcal A}^{i}(X,Y)$ 表示诸群
$\operatorname{Ext}_{\mathcal A}^{0}(X,Y[i])\simeq
\operatorname{Ext}_{\mathcal A}^{0}(X[-i],Y)$（第一章，1.3.8）。这一记号
将在后文得到说明。不过回忆，$D(\mathcal A)$ 的每个特异三角
$(X,X',X'',u,v,w)$（相应地，$(Y,Y',Y'',u,v,w)$）都给出一条无界正合列：
\[
\cdots\longrightarrow \operatorname{Ext}_{\mathcal A}^{0}(X'',Y)
\longrightarrow \operatorname{Ext}_{\mathcal A}^{0}(X',Y)
\longrightarrow \operatorname{Ext}_{\mathcal A}^{0}(X,Y)
\longrightarrow \operatorname{Ext}_{\mathcal A}^{1}(X'',Y)
\longrightarrow\cdots
\]
（相应地，
\[
\cdots\longrightarrow \operatorname{Ext}_{\mathcal A}^{0}(X,Y)
\longrightarrow \operatorname{Ext}_{\mathcal A}^{0}(X,Y')
\longrightarrow \operatorname{Ext}_{\mathcal A}^{0}(X,Y'')
\longrightarrow \operatorname{Ext}_{\mathcal A}^{1}(X,Y)
\longrightarrow\cdots )
\]
（第二章，1.2.1）。

\VerdierSemanticAnchor{ch3d-1-2-6}
\medskip\noindent\textbf{1.2.6.}\quad
设 $M$ 是 $\mathcal A$ 的一个对象集。以 $D_M(\mathcal A)$ 表示
$D(\mathcal A)$ 的满子范畴，它由上同调对象在每个次数都属于 $M$ 的复形
定义。以 $D_M^*(\mathcal A)$ 表示满子范畴
$D_M(\mathcal A)\cap D^*(\mathcal A)$。若由 $M$ 中对象定义的
$\mathcal A$ 的满子范畴是 $\mathcal A$ 的厚子范畴（\textit{即}，它是
阿贝尔范畴、包含函子正合，并且 $M$ 的对象集对扩张封闭），则诸子范畴
$D_M^*(\mathcal A)$ 是 $D(\mathcal A)$ 的满三角子范畴。

\VerdierSemanticAnchor{ch3d-1-2-7}
\medskip\noindent\textbf{1.2.7.}\quad
设 $F\colon\mathcal A\longrightarrow\mathcal B$ 是两个阿贝尔范畴之间的正合
函子，$K(F)\colon K(\mathcal A)\longrightarrow K(\mathcal B)$ 是函子 $F$
在复形同伦范畴上的延拓（第一章，2.5.10）。函子 $K(F)$ 把
$\mathcal A$ 的无环复形变为 $\mathcal B$ 的无环复形，并且是正合的
（第二章，1.3.7），因而由（第二章，2.2.6）定义一个正合函子
$D(F)\colon D(\mathcal A)\longrightarrow D(\mathcal B)$。函子 $D(F)$
把 $D(\mathcal A)$ 的有界（相应地，下有界，相应地，上有界）对象变为
$D(\mathcal B)$ 中同类的对象。在大多数情形下，我们将滥用记号
$F\colon D(\mathcal A)\longrightarrow D(\mathcal B)$ 来表示函子 $D(F)$。

\VerdierSemanticAnchor{ch3d-1-2-8}
\medskip\noindent\textbf{1.2.8.}\quad
典范同构（第二章，1.3.8）
$i\colon K(\mathcal A)^o\longrightarrow K(\mathcal A^o)$ 把
$\operatorname{Ac}(\mathcal A)^o$ 的对象变为
$\operatorname{Ac}(\mathcal A^o)$ 的对象，因而定义了一个仍记作
$i\colon D(\mathcal A)^o\longrightarrow D(\mathcal A^o)$ 的三角范畴同构。
\endgroup
